\section{轨道车辆检修与运维管理认知}

\subsection{地铁车辆段检修流程}
在赛车场车辆段，我们详细了解了地铁车辆检修的工艺流程。如图\ref{fig:metro-maintenance-flow}，首先将列车进行清洗、解编，后再拆分成车体与转向架两部分；转向架再拆分成电机、轮对等部件分别进行维护，后再拼成转向架，最后回到总装车间和车体进行组装。\par{}
\reportFigure{6}
其中，我们主要参观了转向架工区、轮对工区，图\ref{fig:metro-maintenance-scenes} (a)正在进行转向架拆解，图\ref{fig:metro-maintenance-scenes} (b)展示了一系钢弹簧与拆卸下来的轮对，图\ref{fig:metro-maintenance-scenes} (c)是已卸去轮饼的车轴，图\ref{fig:metro-maintenance-scenes} (d)是已经检修完毕、重新安装好轴箱的轮对。需要注意的是，组装工序的步骤也与先前在浦镇公司学习的基本一致，此处不多赘述。\par{}
\reportFigure{7}

\subsection{数字化运维与检测实验室}

\subsubsection{上海地铁数字化运维情况概述}
在梅陇基地，我们了解了上海地铁数字化运维管理中心的构成。其下辖多个机构，包括生产调度中心、车场管理中心、状态监控中心、运输支持中心、应急指挥中心等。\par{}
其中，生产调度中心将生产任务构建成时间一定的工作包，使用自动排程系统（APS）拟定生产计划，促进维修模式转型；车场管理中心运用综合管理系统（DCC）实现作业安全预警、车辆实时定位、机器人安全监测等任务；状态监控中心使用车联网（IOR）系统对车辆运行状态进行实时监控；运输支持中心使用应急处置APP和专家系统等支持地铁系统运作；应急指挥中心则开发应急抢险系统，专门用于应急事态响应工作等。\par{}

\subsubsection{梅陇基地检验检测部实验室概述}
梅陇基地检验检测部实验室于2020年开始筹备建造、2023年正式成立，初衷是通过实验室获取大量的现场检测数据，以满足智能运维的庞大数据需求，其业务包括来样检测、抽样检测、验车检测等，内含齿轮箱润滑油实验室、控制板卡实验室、车辆低压电气实验室、车辆高压电气实验室、全相检测实验室等，下面对其部分实验室主要任务与工作方式进行概述。\par{}
首先是齿轮箱润滑油实验室。齿轮箱润滑油在齿轮箱润滑、缓解齿轮箱磨损上有重要作用，一般需要一年更换一次；然而，一般的4M2T编组按照4节动车、每节车厢2个转向架、每个转向架2个轮对计算，就有$4 \times 2 \times 2 = 16$个齿轮箱。上海地铁全路数千辆车，若每年每车都换油，消耗量巨大，不利于降本增效。因此可以对车辆换下的润滑油进行检测：先取润滑油搅拌均匀后，对水含量、黏度、Cu与Fe等金属离子含量进行检测，会使用到光谱仪。\par{}
接着是控制板卡实验室。该实验室有3个试验台，分别是门控器试验台、继电器试验台和报站系统试验台。门控器试验台就专门用于检测车门控制板卡的信号收发、输出功率、电压波形等；继电器试验台专门对机械式继电器关键部件进行检查，防止关键部件出现机械卡死、热失效等故障；虽然报站与显示不影响列车运行本身，但是报站错误会导致乘客投诉，降低运营效率，因此需要通过报站系统试验台进行检测。\par{}
低压电气实验室有2个试验台，分别是速度传感器试验台与电气元件综合试验台。一般速度传感器采用霍尔式等，根据记录的双脉冲来推算齿轮的转速，进而得知车速；双脉冲比单脉冲的优势就在于其可以同时反映速度大小和方向。速度传感器的准确度主要受其工作环境影响，振动、温度、齿轮箱漏油等可能会使得漏掉脉冲、速度记录出错，因此需要对速度传感器进行检测。而电气元件综合试验台通过模拟振动和温度来对电气元件耐受性进行测试。\par{}
全相检测实验室、车辆高压电气实验室分别主要对走行部、牵引逆变器与受电弓等进行检测，此处笔记较为混乱因此不做论述。\par{}

\subsection{动车组高级修}
上海动车段动车组高级修场场内布局如图\ref{fig:emu-overhaul-yard-layout}所示。其中调试库I与三级检修库也称1号库，中部各种检修区也称2号库，部件检修区与解编修库也称3号库，最下方的三级检修库与调试库（实为整体架车库）也称5号库，而三号库西侧的油漆库等即为4号库；图中的配件中心实为高级修检修控制中心（MCC）。\par{}
\reportFigure{12}

\subsubsection{三级检修库概述}
首先应明确检修库与运用所的区别。检修库指列车定期的检查与修理场所，一般修理较为完全与彻底；运用所则为列车每日运行结束后返回的场所，只进行常规检查，一般以夜班为主，尤其是春运时非常忙碌。\par{}
参观当日（2026年7月16日），中车唐山制造、配属上局南京段的CRH380BL-3545列车正在1号库三级检修库D12股道上进行检修作业，正在进行落成恢复、单车调试、编组的工序，如图\ref{fig:crh380bl-3545-overhaul}所示。D12股道旁的电子屏显示“有电”，说明正在进行通电试验。在通电或断电的过程中，要求附近人员远离车辆与接触网，在断电后还要人为在触网上悬挂接地线。由MCC内的管理界面与国家相关标准\sourcecite{4}可得知，高铁动车组三级检修由接车预检、预检后调车、解编、架车、车体检修、转向架检修、落车、尺寸测量与称重、落成恢复、单车调试、编组、静态调试、动态调试、动调后检查、淋雨试验等工序组成。\par{}
\reportFigure{13}
此外，三级检修库中还有各类设备，大部分设备都是中车青岛四方公司生产的。图\ref{fig:overhaul-equipment}从左至右依次是轴重计算设备、落坑式平台、铁鞋。其中轴重计算设备就是装有压力传感器的地磅式轨道，位于检修库出入口处，主要用于测量列车转向架左右重量是否平衡，一般要求偏差不超过4\%或8\%；落坑式平台可通过下沉的方式将拆卸的转向架从下方转移；铁鞋则是为了防止溜车。\par{}
\reportFigure{14}

\subsubsection{高级修检修控制中心概述}
高级修检修控制中心（MCC）内有专用数字化控制平台，集成了生产指挥、指标管控、数字档案、车间监控等功能，大量运用数字孪生技术，并可通过摄像头实时监视厂房内部情况。数字档案中CRH380BL-3545也在系统中显示，有包括检修流程进度在内的各项详细指标。\par{}
除此外系统还展示了2026一季度短编（应指8节编组）三级检修全路修时排名。修时指从车辆进入厂房到修竣出厂的全过程时长，当前一季度修时中第一名是广州局的12.4天，第二名是上海局的12.7天，第三名是成都局的13.4天……据工人师傅描述，原先进行一次三级检修需要近一个月时间，而当下只需要不到两周，可见检修效率在工人的奋斗与技术的加持下不断提升，真正体现了“科学技术是第一生产力”的道理。\par{}

\subsubsection{部件检修库概述}
3号库中存在一个展览区，通过车上设备制成的艺术品，展示了该部门工作理念——PARTS，即精准（Precision）-能力（Ability）-变革（Revolution）-持续（Timeless）-安全（Security），其中代表A的机器人艺术品如图\ref{fig:parts-artwork}所示。\par{}
\reportFigure{16}
参观区域内主要有两条产线：空调检测和车钩检测产线。其中空调从列车上拆卸后，要经过检修、组装、试验（气密试验、淋雨试验、运转试验、焓差试验等）、交检交验、合格品存放等步骤；而动车组端部的重联用车钩则要经过零部件检查、部件探伤、预装、总装、试验、交检交验、合格品存放等步骤。图\ref{fig:component-overhaul-workstations}则展现了部分工序的工位情况，从左往右依次为：空调检修、空调组装、车钩连挂试验。\par{}
\reportFigure{15}

